\centerline{\bf \Huge gcd+lcm}

\begin{flushright}
        \scriptsize{Эпизод 12. Не заставляй других страдать из-за твоей собственной ненависти, — сказала она}
\end{flushright}


\problem{1}
\textit{Степень вхождения} простого числа $p$ в натуральное $a$ равна $\alpha$, а степень вхождения $p$ в натуральное $b$ равна $\beta$. Докажите, что степень вхождения $p$ в НОД( $a, b$ ) равна $\min (\alpha, \beta)$, а степень вхождения $p$ в $\operatorname{HOK}(a, b)$ равна $\max (\alpha, \beta)$.

\problem{2}
Докажите, что $\operatorname{HOД}(a, b) \cdot \operatorname{HOK}(a, b)=a b$.

\problem{3}
Два натуральных числа $a$ и $b$ друг на друга не делятся, при этом НОД( $a, b)=50$, а $\operatorname{HOK}(a, b)=1000$. Найдите эти числа.

\problem{4}
Пусть $a, b, c$ - натуральные числа. Могут ли наибольшие общие делители пар чисел $a$ и $b, b$ и $c, c$ и $a$ равняться $30!+111,40!+234$ и $50!+666$ соответственно?

\problem{5}
Сколько существует пар натуральных чисел $(a, b)$, у которых $\operatorname{HOK}(a, b)=2^3 \cdot 3^2 \cdot 5 \cdot 7^2=17640$, а НОД $(a, b)=12$ ? (пары неупорядоченные, то есть $(a, b)$ и $(b, a)$ считайте одинаковыми)

\problem{6}
Можно ли вместо звездочек вставить в выражение

$$
\operatorname{HOK}(*, *, *)-\operatorname{HOK}(*, *, *)=2009
$$


в некотором порядке шесть последовательных натуральных чисел так, чтобы равенство стало верным?